% =============================================================================
% Template — resolução de listas (MAC0425 ou outras disciplinas IME)
% Compilação: pdflatex (recomendado) ou latexmk -pdf template-lista.tex
% =============================================================================
\documentclass[a4paper,11pt]{article}

\usepackage[utf8]{inputenc}
\usepackage[T1]{fontenc}
\usepackage[brazil]{babel}

\usepackage[margin=2.4cm, top=2.7cm, bottom=2.6cm, headheight=14pt]{geometry}
\usepackage{microtype}
\usepackage{xcolor}
\usepackage{enumitem}
\usepackage{amsmath, amssymb, amsthm}
\usepackage{mathtools}
\usepackage{graphicx}
\usepackage{booktabs}
\usepackage{hyperref}
\usepackage{bookmark}
\usepackage{etoolbox}
\usepackage{fancyhdr}
\usepackage{lastpage}
\usepackage{tikz}
\usetikzlibrary{trees, positioning, arrows.meta, shadows, calc}

\usepackage{parskip}
\usepackage[linesnumbered,ruled,vlined,portuguese]{algorithm2e}
\SetKwProg{Fn}{Função}{:}{Fim}
\usepackage{listings}
\usepackage[most]{tcolorbox}

% --- Hiperlinks discretos ---
\hypersetup{
  colorlinks=true,
  linkcolor=IMEblue!85!black,
  urlcolor=IMEblue!75!black,
  citecolor=IMEblue!75!black,
}

% --- Paleta ---
\definecolor{IMEblue}{RGB}{0, 72, 125}
\definecolor{IMEgray}{RGB}{88, 89, 91}
\definecolor{IMElightgray}{RGB}{245, 247, 250}

% --- Configuração do Listings (Estilo para Códigos) ---
\lstset{
  basicstyle=\ttfamily\small,
  keywordstyle=\color{IMEblue}\bfseries,
  commentstyle=\color{IMEgray}\itshape,
  stringstyle=\color{green!50!black},
  backgroundcolor=\color{IMElightgray},
  frame=single,
  rulecolor=\color{IMEgray!40},
  breaklines=true,
  showstringspaces=false
}

% --- Atalhos Matemáticos ---
\DeclareMathOperator*{\argmax}{arg\,max}
\DeclareMathOperator*{\argmin}{arg\,min}
\DeclareMathOperator{\E}{\mathbb{E}}
\DeclareMathOperator{\Var}{Var}
\DeclareMathOperator{\Cov}{Cov}
\newcommand{\mbf}[1]{\mathbf{#1}}

% ----======== PREENCHA  ========----
\newcommand{\CodigoDisciplina}{MAC0105}
\newcommand{\NomeDisciplina}{Fundamentos de Matemática para a Computação}
\newcommand{\NumeroLista}{2}
\newcommand{\TituloListaExtra}{Lista de exercícios 2}
\newcommand{\NomeAluno}{Leonardo Heidi Almeida Murakami}
\newcommand{\NUSP}{11260186}
\newcommand{\DataEntrega}{\today}
% ----=========================================----

% Cabeçalho e rodapé
\pagestyle{fancy}
\fancyhf{}
\fancyhead[L]{\small\color{IMEgray}\NomeAluno}
\fancyhead[R]{\small\color{IMEgray}NUSP \NUSP}
\renewcommand{\headrulewidth}{0.4pt}
\renewcommand{\headrule}{\hbox to\headwidth{\color{IMEgray!40}\leaders\hrule height \headrulewidth\hfill}}
\fancyfoot[L]{\small\color{IMEgray}\CodigoDisciplina\ -- Lista \NumeroLista}
\fancyfoot[C]{\small\color{IMEgray}\NomeDisciplina}
\fancyfoot[R]{\small\color{IMEgray}\thepage\ / \pageref{LastPage}}

\fancypagestyle{plain}{
  \fancyhf{}
  \fancyfoot[R]{\small\color{IMEgray}\thepage\ / \pageref{LastPage}}
  \renewcommand{\headrulewidth}{0pt}
}

% Teoremas / demonstrações
\theoremstyle{definition}
\newtheorem{definicao}{Definição}[section]

\theoremstyle{plain}
\newtheorem{teorema}{Teorema}[section]

\theoremstyle{remark}
\newtheorem*{observacao}{Observação}

% Ambiente: enunciado
\newenvironment{enunciado}
  {\par\noindent\color{IMEgray}\itshape}
  {\par\medskip}

% --- EXERCICIOS ---

% 1. Exercício com contador automático
\newcounter{questao}[section]
\newtcolorbox[use counter=questao]{questao}[1][]{
  enhanced,
  colback=white,
  colframe=IMElightgray,
  fonttitle=\bfseries\color{IMEblue},
  coltitle=IMEblue,
  colbacktitle=IMElightgray,
  title={Questão~\thetcbcounter\ifstrempty{#1}{}{ \textnormal{\color{black}#1}}},
  boxrule=1pt,
  arc=2pt,
  outer arc=2pt,
  top=3mm,
  bottom=3mm,
  left=3mm,
  right=3mm,
  breakable
}

\newtcolorbox{questaoManually}[1]{%
  enhanced,
  colback=white,
  colframe=IMElightgray,
  fonttitle=\bfseries\color{IMEblue},
  coltitle=IMEblue,
  colbacktitle=IMElightgray,
  title={#1},
  boxrule=1pt,
  arc=2pt,
  outer arc=2pt,
  top=3mm,
  bottom=3mm,
  left=3mm,
  right=3mm,
  breakable
}

\newcommand{\makelistatitle}{%
  \thispagestyle{plain}
  \noindent
  \begin{minipage}[t]{0.70\linewidth}
    {\LARGE\bfseries\color{IMEblue}\CodigoDisciplina}\par\vspace{0.15em}
    {\large\color{IMEgray}\NomeDisciplina}\par\vspace{1.1em}
    {\Large\bfseries Lista \NumeroLista
      \ifdefempty{\TituloListaExtra}{}{ --- \TituloListaExtra}%
    }
  \end{minipage}\hfill
  \begin{minipage}[t]{0.28\linewidth}
    \raggedleft\small\color{IMEgray}
    \textbf{\NomeAluno}\\[0.25em]
    NUSP \NUSP\\[0.5em]
    \DataEntrega
  \end{minipage}
  \par\vspace{1.2em}
  \noindent{\color{IMEgray!35}\rule{\linewidth}{0.6pt}}
  \par\vspace{1em}
}

\begin{document}

\makelistatitle

\begin{questaoManually}{Exercício 1}
Determine se as expressões lógicas são tautologias.

\begin{enumerate}
\item \textbf{$((p \to q) \land (q \to r)) \to (p \to r)$} \\
\textit{Resposta: Tautologia.} \\
\textit{Justificativa:} Silogismo hipotético. Por transformações equivalentes:
\begin{align*}
((p \to q) \land (q \to r)) \to (p \to r) &\iff \neg ( (\neg p \lor q) \land (\neg q \lor r) ) \lor (\neg p \lor r) \\
&\iff \neg (\neg p \lor q) \lor \neg (\neg q \lor r) \lor (\neg p \lor r) \tag{De Morgan}\\
&\iff (p \land \neg q) \lor (q \land \neg r) \lor \neg p \lor r \\
&\iff ((p \land \neg q) \lor \neg p) \lor ((q \land \neg r) \lor r) \tag{Associatividade}\\
&\iff (\neg p \lor \neg q) \lor (q \lor r) \tag{Absorção}\\
&\iff \neg p \lor (\neg q \lor q) \lor r \\
&\iff \neg p \lor V \lor r \iff V
\end{align*}

\item \textbf{$((p \to q) \land (p \to \neg q)) \to \neg p$} \\
\textit{Resposta: Tautologia.} \\
\textit{Justificativa:}
\begin{align*}
((p \to q) \land (p \to \neg q)) &\iff (\neg p \lor q) \land (\neg p \lor \neg q) \\
&\iff \neg p \lor (q \land \neg q) \tag{Distributividade}\\
&\iff \neg p \lor F \iff \neg p
\end{align*}
A expressão inteira se torna $\neg p \to \neg p \iff \neg(\neg p) \lor \neg p \iff p \lor \neg p \iff V$.

\item \textbf{$(p \land (p \to q)) \to q$} \\
\textit{Resposta: Tautologia.} \\
\textit{Justificativa:} Conhecida como \textit{Modus Ponens}.
\begin{align*}
(p \land (p \to q)) \to q &\iff (p \land (\neg p \lor q)) \to q \\
&\iff ((p \land \neg p) \lor (p \land q)) \to q \\
&\iff (F \lor (p \land q)) \to q \iff (p \land q) \to q \\
&\iff \neg(p \land q) \lor q \iff \neg p \lor \neg q \lor q \\
&\iff \neg p \lor V \iff V
\end{align*}

\item \textbf{$(p \to q) \lor (q \to p)$} \\
\textit{Resposta: Tautologia.} \\
\textit{Justificativa:}
\begin{align*}
(p \to q) \lor (q \to p) &\iff (\neg p \lor q) \lor (\neg q \lor p) \\
&\iff (\neg p \lor p) \lor (\neg q \lor q) \\
&\iff V \lor V \iff V
\end{align*}

\item \textbf{$(p \oplus q) \leftrightarrow \neg(p \leftrightarrow q)$} \\
\textit{Resposta: Tautologia.} \\
\textit{Justificativa:} Sabemos que $p \leftrightarrow q \iff (p \land q) \lor (\neg p \land \neg q)$. Negando a bicondicional:
\begin{align*}
\neg(p \leftrightarrow q) &\iff \neg((p \land q) \lor (\neg p \land \neg q)) \\
&\iff \neg(p \land q) \land \neg(\neg p \land \neg q) \\
&\iff (\neg p \lor \neg q) \land (p \lor q) \\
&\iff (\neg p \land p) \lor (\neg p \land q) \lor (\neg q \land p) \lor (\neg q \land q) \\
&\iff F \lor (\neg p \land q) \lor (p \land \neg q) \lor F \\
&\iff (p \land \neg q) \lor (\neg p \land q)
\end{align*}
A última linha é a exata definição de $p \oplus q$.

\item \textbf{$(p \land q) \lor (\neg p \land \neg q)$} \\
\textit{Resposta: Não é tautologia.} \\
\textit{Justificativa:} Esta é a definição de $p \leftrightarrow q$. A expressão é falsa se $p$ e $q$ tiverem valores diferentes. Exemplo: se $p = V$ e $q = F$, $(V \land F) \lor (F \land V) \iff F \lor F \iff F$.

\item \textbf{$p \land (q \lor r) \leftrightarrow (p \land q) \lor (p \land r)$} \\
\textit{Resposta: Tautologia.} \\
\textit{Justificativa:} Trata-se da Lei Distributiva da conjunção sobre a disjunção, que é uma equivalência lógica fundamental.

\item \textbf{$((p \to q) \land (q \to p)) \to (p \leftrightarrow q)$} \\
\textit{Resposta: Tautologia.} \\
\textit{Justificativa:} A premissa $((p \to q) \land (q \to p))$ é exatamente a definição da bicondicional $(p \leftrightarrow q)$. Logo, temos $(p \leftrightarrow q) \to (p \leftrightarrow q)$, o que é trivialmente $V$.

\item \textbf{$((p \to q) \land (p \to \neg q)) \to p$} \\
\textit{Resposta: Não é tautologia.} \\
\textit{Justificativa:} Vimos no item 2 que a premissa é equivalente a $\neg p$. A expressão reduz-se a $\neg p \to p \iff \neg(\neg p) \lor p \iff p \lor p \iff p$. Como o resultado depende do valor de $p$, não é uma tautologia (falsa quando $p = F$).
\end{enumerate}

\end{questaoManually}
\newpage \begin{questaoManually}{Exercício 2}
Com base na imagem do tabuleiro fornecida:

\begin{enumerate}
\item \textbf{$\forall x, y (Triângulo(x, y) \to \neg Azul(x, y))$} \\
\textit{Resposta: Verdadeiro.} \\
\textit{Justificativa:} A afirmação diz que ``todo triângulo não é azul'' . Inspecionando o tabuleiro, os únicos triângulos presentes são verdes (nas casas $(1, 8)$ e $(3, 6)$) e vermelhos (nas casas $(2, 2)$ e $(6, 4)$) . Como não existe nenhum triângulo azul na imagem, a condição lógica é totalmente satisfeita [1, 2].

\item \textbf{$\forall x, y (\neg Azul(x, y) \to Triângulo(x, y))$} \\
\textit{Resposta: Falso.} \\
\textit{Justificativa:} A proposição dita que ``qualquer posição que não possua uma peça azul obrigatoriamente contém um triângulo'' . Podemos facilmente encontrar contraexemplos que invalidam isso, como a casa $(2, 3)$, que possui um quadrado verde (não é azul e também não é triângulo) . 

\item \textbf{$\exists x, y (\neg Azul(x, y) \to (Triângulo(x, y) \lor Vazio(x, y)))$} \\
\textit{Resposta: Verdadeiro.} \\
\textit{Justificativa:} A afirmação existencial pede que encontremos pelo menos uma casa onde a implicação seja válida . Tomando a casa $(1, 8)$ como exemplo, temos um triângulo verde . A premissa $\neg Azul(1, 8)$ é verdadeira, e a conclusão de ser um triângulo ou estar vazio também é verdadeira, confirmando a proposição como um todo [1, 2].

\item \textbf{$\exists x, y (\neg Vermelho(x, y) \land Triângulo(x, y))$} \\
\textit{Resposta: Verdadeiro.} \\
\textit{Justificativa:} A proposição afirma que ``existe ao menos um triângulo que não é vermelho'' . Os triângulos de cor verde localizados nas coordenadas $(1, 8)$ e $(3, 6)$ satisfazem exatamente esta exigência .

\item \textbf{$\forall x \exists y (\neg Vazio(x, y) \land \neg Triângulo(x, y))$} \\
\textit{Resposta: Verdadeiro.} \\
\textit{Justificativa:} Afirma-se que ``em cada uma das 8 colunas ($x$), existe ao menos uma linha ($y$) contendo uma peça que não é um triângulo'' . Inspecionando o tabuleiro da esquerda para a direita , encontramos as seguintes peças não triangulares em cada coluna respectiva:
Círculo azul em $(1, 4)$; Quadrado verde em $(2, 3)$; Quadrado vermelho em $(3, 2)$; Quadrado azul em $(4, 6)$; Círculo azul em $(5, 2)$; Círculo vermelho em $(6, 3)$; Quadrado vermelho em $(7, 1)$; Círculo azul em $(8, 1)$. Logo, todas as colunas cumprem o requisito .

\item \textbf{$\exists y \forall x (Vazio(x, y) \lor (\neg Triângulo(x, y) \land Verde(x, y)))$} \\
\textit{Resposta: Verdadeiro.} \\
\textit{Justificativa:} Esta proposição exige que exista uma linha inteira ($y$) formada unicamente por casas vazias ou por peças verdes que não sejam triângulos . Observando a \textbf{linha 5}, notamos que as posições ocupadas são $(5, 5)$ e $(7, 5)$, as quais contêm, respectivamente, um círculo verde e um quadrado verde . Todas as demais casas desta linha estão estritamente vazias, o que torna a proposição verdadeira [1, 2].

\item \textbf{$\exists x \forall y (Círculo(x, y) \lor \neg(Quadrado(x, y) \lor Triângulo(x, y)))$} \\
\textit{Resposta: Verdadeiro.} \\
\textit{Justificativa:} Aplicando a Lei de De Morgan no trecho negado, temos a equivalência $\exists x \forall y (Círculo(x, y) \lor (\neg Quadrado(x, y) \land \neg Triângulo(x, y)))$ . Lendo em português, isso significa que ``existe uma coluna ($x$) inteira onde todas as peças presentes são círculos ou a casa está vazia'' (ou seja, livre de quadrados e triângulos) . As \textbf{colunas 5 e 8} atendem perfeitamente a isso: a coluna 5 possui apenas um círculo azul em $(5, 2)$ e um círculo verde em $(5, 5)$, enquanto a coluna 8 possui apenas um círculo azul na ponta em $(8, 1)$ . Nenhuma delas possui triângulos ou quadrados .
\end{enumerate}

\end{questaoManually}
\newpage \begin{questaoManually}{Exercício 3}
Com base nos predicados descritos, transcrevemos as afirmações para a linguagem lógica:

\begin{enumerate}
\item \textbf{Todo círculo azul contém um quadrado a esquerda.}
\begin{equation}
\forall c (Círculo(c) \land Azul(c) \to \exists d (Quadrado(d) \land Esquerda(d, c)))
\end{equation}

\item \textbf{Todo objeto azul contém um objeto verde abaixo.}
\begin{equation}
\forall c (Azul(c) \to \exists d (Verde(d) \land Acima(c, d)))
\end{equation}

\item \textbf{Existe um quadrado com um triângulo abaixo.}
\begin{equation}
\exists c \exists d (Quadrado(c) \land Triângulo(d) \land Acima(c, d))
\end{equation}

\item \textbf{Existe uma casa vazia com um quadrado acima e um triângulo a esquerda.}
\begin{equation}
\exists c, d, e (Vazio(c) \land Quadrado(d) \land Acima(d, c) \land Triângulo(e) \land Esquerda(e, c))
\end{equation}

\item \textbf{Existe uma casa com um objeto verde nos dois lados.}
\begin{equation}
\exists c, d, e (Verde(d) \land Esquerda(d, c) \land Verde(e) \land Esquerda(c, e))
\end{equation}

\item \textbf{Existe uma coluna contendo apenas um círculo na sua casa mais baixa.}
\begin{equation}
\exists c (Círculo(c) \land \neg \exists d (Acima(c, d)) \land \forall e (Acima(e, c) \to \neg Círculo(e)))
\end{equation}
\end{enumerate}

\end{questaoManually}
\newpage \begin{questaoManually}{Exercício 4}

\textbf{(a) Se $x$ é ímpar, então $x^2$ é ímpar.}

\textit{Demonstração (Prova Direta):}
Seja $x$ um número inteiro ímpar. Pela definição, podemos escrever $x = 2k + 1$, para algum $k \in \mathbb{Z}$.
Calculando o valor de $x^2$:
\begin{align*}
x^2 &= (2k + 1)^2 \\
&= 4k^2 + 4k + 1 \\
&= 2(2k^2 + 2k) + 1
\end{align*}
Como $k \in \mathbb{Z}$, temos que $2k^2 + 2k$ também é um número inteiro. Seja $m = 2k^2 + 2k$. Podemos reescrever:
\begin{equation}
x^2 = 2m + 1
\end{equation}
Esta é precisamente a forma de um número ímpar, concluindo a prova. \blacksquare

\textbf{(b) Suponha que $x$ e $y$ são números reais. Se $y^3 + yx^2 \le x^3 + xy^2$, então $y \le x$.}

\textit{Demonstração (Contrapositiva):}
Provaremos a contrapositiva: Se $y > x$, então $y^3 + yx^2 > x^3 + xy^2$.
Suponha $y > x$, o que nos dá $y - x > 0$.
Vamos analisar a expressão $(y^3 + yx^2) - (x^3 + xy^2)$:
\begin{align*}
y^3 - x^3 + yx^2 - xy^2 &= y^3 - xy^2 + yx^2 - x^3 \\
&= y^2(y - x) + x^2(y - x) \\
&= (y^2 + x^2)(y - x)
\end{align*}
Como impomos $y > x$, é impossível que $y$ e $x$ sejam simultaneamente nulos. Assim, temos sempre $y^2 + x^2 > 0$.
Sabendo por hipótese que $(y - x) > 0$, concluímos que o produto de ambos os fatores deve ser positivo:
\begin{align*}
(y^2 + x^2)(y - x) &> 0 \\
y^3 + yx^2 - x^3 - xy^2 &> 0 \implies y^3 + yx^2 > x^3 + xy^2
\end{align*}
Como a contrapositiva é verdadeira, a afirmação original é provada. \blacksquare

\end{questaoManually}
\newpage \begin{questaoManually}{Exercício 5}

\textbf{(a) $x < n \iff \lfloor x \rfloor < n$}

\textit{Resposta: Verdadeiro.}

\textit{Demonstração:}
($\implies$) Se $x < n$, por definição do piso sabemos que $\lfloor x \rfloor \le x$. Por transitividade, $\lfloor x \rfloor < n$.
($\impliedby$) Se $\lfloor x \rfloor < n$, como $\lfloor x \rfloor$ e $n$ são números inteiros, a desigualdade estrita garante que $\lfloor x \rfloor \le n - 1$.
Pela definição de piso, temos sempre a propriedade fundamental $x < \lfloor x \rfloor + 1$. Substituindo a limitação que acabamos de deduzir:
\begin{equation}
x < \lfloor x \rfloor + 1 \le (n - 1) + 1 = n \implies x < n \ \blacksquare
\end{equation}


\textbf{(b) $x \ge n \iff \lceil x \rceil < n$}

\textit{Resposta: Falso.}

\textit{Contraexemplo:}
Considere $x = 1.5$ e o inteiro $n = 1$.
Avaliando a premissa à esquerda, temos $1.5 \ge 1$, o que é \textbf{Verdadeiro}.
No entanto, $\lceil 1.5 \rceil = 2$. Verificando o lado direito, teríamos $2 < 1$, o que é \textbf{Falso}.
Como a condicional $(V \iff F)$ resulta em falso, demonstramos que a propriedade não é válida em todos os casos. \blacksquare


\end{questaoManually}
\newpage \begin{questaoManually}{Exercício 6}
Prove que $H_{2^n} \ge 1 + n/2$ para todo inteiro $n \ge 0$.

\textit{Demonstração (Por Indução):}
Sabe-se que o número harmônico é definido por $H_m = \sum_{i=1}^m \frac{1}{i}$.

\textbf{Passo Base ($n=0$):}
\begin{equation}
H_{2^0} = H_1 = \sum_{i=1}^1 \frac{1}{i} = 1
\end{equation}
Verificando do outro lado da desigualdade: $1 + \frac{0}{2} = 1$.
Como $1 \ge 1$ é verdadeiro, a base está correta.

\textbf{Passo Indutivo:}
Suponha que a propriedade seja verdadeira para um inteiro genérico $k \ge 0$, o que nos dá nossa \textit{Hipótese de Indução (H.I.)}:
\begin{equation}
H_{2^k} \ge 1 + \frac{k}{2}
\end{equation}
Queremos demonstrar que o caso $k+1$ também é válido: $H_{2^{k+1}} \ge 1 + \frac{k+1}{2}$.
Expandindo a soma para a potência $2^{k+1}$:
\begin{align*}
H_{2^{k+1}} &= \sum_{i=1}^{2^{k+1}} \frac{1}{i} \\
&= \sum_{i=1}^{2^k} \frac{1}{i} + \sum_{i=2^k+1}^{2^{k+1}} \frac{1}{i} \\
&= H_{2^k} + \sum_{i=2^k+1}^{2^{k+1}} \frac{1}{i}
\end{align*}
A segunda somatória possui exatamente $2^{k+1} - 2^k = 2^k(2 - 1) = 2^k$ termos. O menor valor contido nela encontra-se no final (quando $i = 2^{k+1}$). Portanto, ao substituir todas as frações por $\frac{1}{2^{k+1}}$, montamos um limitante inferior:
\begin{equation}
\sum_{i=2^k+1}^{2^{k+1}} \frac{1}{i} \ge \sum_{i=2^k+1}^{2^{k+1}} \frac{1}{2^{k+1}} = 2^k \left( \frac{1}{2^{k+1}} \right) = \frac{1}{2}
\end{equation}
Voltando à nossa equação e substituindo a H.I. e o limitante calculado:
\begin{align*}
H_{2^{k+1}} &\ge H_{2^k} + \frac{1}{2} \\
H_{2^{k+1}} &\ge \left( 1 + \frac{k}{2} \right) + \frac{1}{2} \\
H_{2^{k+1}} &\ge 1 + \frac{k+1}{2}
\end{align*}
O passo indutivo foi garantido. Pela indução finita, a relação $H_{2^n} \ge 1 + n/2$ se sustenta para todo inteiro $n \ge 0$. \blacksquare

\end{questaoManually}
\newpage \begin{questaoManually}{Exercício 7}
Mostre que para todo inteiro positivo $n$ valem os seguinte itens .

\textbf{(a)} $\sum_{i=1}^n i^3 = \left( \sum_{i=1}^n i \right)^2$

\textit{Demonstração (Por Indução):}
Sabemos que a soma dos $n$ primeiros inteiros é dada por $\sum_{i=1}^n i = \frac{n(n+1)}{2}$. Logo, queremos provar a equivalência $\sum_{i=1}^n i^3 = \frac{n^2(n+1)^2}{4}$ .

\textbf{Passo Base ($n=1$):}
O lado esquerdo resulta em $1^3 = 1$. O lado direito é $\frac{1^2(1+1)^2}{4} = \frac{4}{4} = 1$. A base é verdadeira.

\textbf{Passo Indutivo:}
Suponha que a propriedade seja verdadeira para um inteiro $k \ge 1$, que será nossa Hipótese de Indução (H.I.):
\begin{equation}
\sum_{i=1}^k i^3 = \frac{k^2(k+1)^2}{4}
\end{equation}
Queremos provar para o sucessor $n = k+1$:
\begin{align*}
\sum_{i=1}^{k+1} i^3 &= \sum_{i=1}^k i^3 + (k+1)^3 \\
&= \frac{k^2(k+1)^2}{4} + (k+1)^3 \tag{Substituição da H.I.} \\
&= \frac{k^2(k+1)^2 + 4(k+1)^3}{4} \\
&= \frac{(k+1)^2 [k^2 + 4(k+1)]}{4} \tag{Colocando $(k+1)^2$ em evidência} \\
&= \frac{(k+1)^2 (k^2 + 4k + 4)}{4} \\
&= \frac{(k+1)^2 (k+2)^2}{4}
\end{align*}
O resultado alcançado é exatamente a expressão correspondente a $n = k+1$. Por indução, a igualdade é válida para todo inteiro positivo. \blacksquare
\bigskip

\textbf{(b)} $\sum_{i=1}^n i^2 = \frac{2n^3 + 3n^2 + n}{6}$

\textit{Demonstração (Por Indução):}
Note inicialmente que o numerador do lado direito pode ser fatorado como $n(2n^2 + 3n + 1) = n(n+1)(2n+1)$ . 

\textbf{Passo Base ($n=1$):}
Esquerda: $1^2 = 1$. Direita: $\frac{2(1)^3 + 3(1)^2 + 1}{6} = \frac{6}{6} = 1$. A base é verdadeira.

\textbf{Passo Indutivo:}
Assuma que para um $k \ge 1$ valha $\sum_{i=1}^k i^2 = \frac{2k^3 + 3k^2 + k}{6}$.
Avaliando o caso para $n = k+1$:
\begin{align*}
\sum_{i=1}^{k+1} i^2 &= \sum_{i=1}^k i^2 + (k+1)^2 \\
&= \frac{2k^3 + 3k^2 + k}{6} + (k+1)^2 \tag{H.I.}\\
&= \frac{2k^3 + 3k^2 + k + 6(k^2 + 2k + 1)}{6} \\
&= \frac{2k^3 + 9k^2 + 13k + 6}{6}
\end{align*}
Agora, vamos verificar o resultado da fórmula do polinômio original aplicada diretamente a $(k+1)$:
\begin{align*}
\frac{2(k+1)^3 + 3(k+1)^2 + (k+1)}{6} &= \frac{2(k^3+3k^2+3k+1) + 3(k^2+2k+1) + k + 1}{6} \\
&= \frac{2k^3 + 6k^2 + 6k + 2 + 3k^2 + 6k + 3 + k + 1}{6} \\
&= \frac{2k^3 + 9k^2 + 13k + 6}{6}
\end{align*}
Como ambas as expansões convergem exatamente ao mesmo polinômio, o passo indutivo é válido. \blacksquare

\end{questaoManually}
\newpage \begin{questaoManually}{Exercício 8}
Prove que para um inteiro $n \ge 10$ temos que $100n \le 2^n$ .

\textit{Demonstração (Por Indução):}

\textbf{Passo Base ($n=10$):}
Calculando os valores: $100(10) = 1000$ e $2^{10} = 1024$. Como a afirmação $1000 \le 1024$ é verdadeira, a base é garantida .

\textbf{Passo Indutivo:}
Suponha que a desigualdade se sustente para um inteiro $k \ge 10$, isto é, a H.I. é $100k \le 2^k$.
Deseja-se provar para $k+1$, isto é, $100(k+1) \le 2^{k+1}$. Analisando o lado esquerdo:
\begin{equation}
100(k+1) = 100k + 100
\end{equation}
Sabemos que $k \ge 10$. Portanto, é seguro afirmar que $k \ge 1$. Multiplicando por 100 em ambos os lados, concluímos que $100k \ge 100$. Substituindo esse limitante na nossa equação:
\begin{equation}
100k + 100 \le 100k + 100k = 2(100k)
\end{equation}
Pela nossa Hipótese de Indução, $100k \le 2^k$. Usando essa limitação na equação anterior:
\begin{equation}
2(100k) \le 2(2^k) = 2^{k+1}
\end{equation}
Pela propriedade da transitividade das desigualdades, resulta que $100(k+1) \le 2^{k+1}$. A relação está demonstrada. \blacksquare

\end{questaoManually}
\newpage \begin{questaoManually}{Exercício 9}
Prove por indução que todo polígono convexo com $n$ lados pode ser dividido em triângulos usando $n - 3$ diagonais .

\textit{Demonstração (Por Indução):}
Seja $n$ o número de lados do polígono convexo, com $n \ge 3$.

\textbf{Passo Base ($n=3$):}
Um polígono convexo de 3 lados é, por definição, um triângulo . Ele não possui diagonais internas, necessitando de $3 - 3 = 0$ diagonais para ser "triangulado". A proposição atende ao caso base.

\textbf{Passo Indutivo:}
Nossa hipótese (H.I.) dita que todo polígono convexo de $k$ lados ($k \ge 3$) pode ser fatiado perfeitamente em triângulos com o auxílio de $k - 3$ diagonais.
Tomemos agora um polígono convexo de $k+1$ lados. Escolhemos livremente três vértices sucessivos do seu perímetro, os quais chamaremos de $V_1, V_2$ e $V_3$. Traçamos uma diagonal conectando diretamente $V_1$ a $V_3$. 
Ao criar essa conexão, isolamos o triângulo $\triangle V_1 V_2 V_3$ do restante da figura (pelo fato do polígono ser estritamente convexo, a diagonal $V_1V_3$ sempre cruzará o interior do polígono, tornando este corte válido).
Analisando a estrutura após o corte:
\begin{itemize}
    \item Adicionamos \textbf{1} diagonal.
    \item O vértice $V_2$ perdeu conexão geométrica com o polígono restante, transformando nossa figura principal num polígono com vértices $V_1, V_3, V_4, \dots, V_{k+1}$. Ele é composto por $(k+1) - 1 = k$ lados.
\end{itemize}
Pela nossa H.I., o polígono restante com $k$ lados precisará usar $k - 3$ diagonais para ser completamente coberto por triângulos.
Somando as diagonais do processo completo, teremos a diagonal do corte do $V_2$ unida às utilizadas pela figura menor:
\begin{equation}
Total = 1 + (k - 3) = (k + 1) - 3
\end{equation}
A relação $(k+1)-3$ foi atingida, garantindo assim a veracidade universal da tese pelo Princípio da Indução. \blacksquare

\end{questaoManually}
\newpage \begin{questaoManually}{Exercício 10}
Considere uma fila formada por gatos e capivaras. Prove que se o primeiro animal da fila é uma capivara e o último é um gato, então existe uma capivara imediatamente na frente de algum gato .

\textit{Demonstração (Por Contradição):}
Suponha, por absurdo, que a afirmação seja falsa. Ou seja, admita como premissa verdadeira que \textbf{nenhuma} capivara na fila se encontra na frente direta de um gato .
Isto força uma regra de encadeamento estrita: se o animal alocado na posição $i$ da fila for uma capivara, obrigatoriamente o animal atrás dele na posição $i+1$ não poderá ser um gato. A única opção restante (desconsiderando o fim da fila) é que o indivíduo $i+1$ seja outra capivara.
É dado em enunciado que o primeiro animal (posição 1) é uma capivara .
Pela nossa suposição imposta acima, sendo o animal 1 uma capivara, o animal 2 tem que ser uma capivara. Por repetição lógica, se o animal 2 é uma capivara, o 3 é uma capivara, e assim o comportamento se estende recursivamente até o fim da fila, obrigando todos os indivíduos da fila a pertencerem à mesma espécie.
Isto culmina numa consequência inevitável de que o último animal na fila é uma capivara. Entretanto, o enunciado relata explicitamente que a última posição é ocupada por um \textbf{gato} , gerando um evidente Paradoxo/Contradição.
Ao invalidarmos a suposição original, confirma-se o seu oposto: é matematicamente certo que deve existir no mínimo uma posição em que uma capivara precede um gato na fila. \blacksquare

\end{questaoManually}
\newpage \begin{questaoManually}{Exercício 11}
Prove que, para qualquer modo que você quebrar a barra (formada por $n$ quadradinhos) de acordo com a regra, você sempre precisará de exatamente $n - 1$ quebras para separar todos os quadradinhos .

\textit{Demonstração (Por Indução Forte):}
Seja $P(n)$ a afirmação: "Uma barra composta de $n$ blocos necessitará de exatamente $n-1$ fraturas retas para ser convertida em peças unitárias" .

\textbf{Passo Base ($n=1$):}
Uma peça contendo apenas 1 bloco já se encontra perfeitamente separada na dimensão unitária, exigindo $1 - 1 = 0$ quebras . A declaração base é correta.

\textbf{Passo Indutivo (Indução Forte):}
Consideremos a afirmação $P(i)$ verdadeira para todos os tamanhos inteiros $i$ no intervalo de $1 \le i < k$.
Possuímos agora uma barra não quebrada com $k$ quadradinhos totais ($k \ge 2$). Efetuaremos nossa primeira quebra reta nas ranhuras. Independentemente da ranhura escolhida, essa atitude particionará a nossa barra principal de tamanho $k$ em dois pequenos segmentos avulsos de tamanhos $x$ e $y$.
É garantido que a união de ambos dá a barra original ($x + y = k$) e que ambos contêm ao menos 1 quadrado ($x \ge 1$ e $y \ge 1$). Isso garante que os valores $x$ e $y$ são estritamente menores que $k$, credenciando-os a passar pela nossa Hipótese de Indução Forte:
\begin{itemize}
    \item Para reduzir o trecho menor com tamanho $x$ a unidades individuais, necessitamos de exatas $x - 1$ quebras.
    \item Para o segmento $y$, análogo, necessitaremos de exatas $y - 1$ quebras.
\end{itemize}
Contando a quebra primária no total, a quantidade global de ranhuras rompidas será:
\begin{equation}
Total = 1 + (x - 1) + (y - 1) = x + y - 1
\end{equation}
Sabendo que a contagem inteira $x+y$ equivale a $k$, o resultado se transforma em $k - 1$.
Demonstramos assim que, alheio a escolha da ranhura e dimensão lateral, as partições se reduzem a $n-1$ quebras no caso geral, cumprindo o passo da indução. \blacksquare
\end{questaoManually}
\end{document}